\section{One actual original vertical Gauss point for the whole component}

Let $k$ be a field, $Y$ integral, $f:Y\to\operatorname{Spec}k[t]$ proper,
and $i:C\hookrightarrow Y$ a genuine closed immersion with $C$ integral.
Write $\eta=i(\eta_C)$. At an original $x\in C$ put
$p=i(x)$ and $S=\mathcal O_{Y,p}$. Assume this actual $S$ is Noetherian,
and an original prime element $\pi\in S$ satisfies
$\ker(i_x^\#)=(\pi)$. Retain the actual original coefficient map of $f$,
its parameter $t_p\ne0$, and $\pi\mid t_p$. No field isomorphism,
valuation family, center, normalized norm, or native-point equality is input.

\begin{lemma}[The actual vertical DVR morphism has the actual component center]
\label{lem:gauss-common-original-vertical-dvr-morphism}
\uses{thm:gauss-common-original-component-stalk,def:gauss-surface-original-residual-vertical-rings,def:gauss-surface-original-vertical-fraction-map,lem:gauss-surface-original-vertical-fraction-identities,lem:actual-specified-field-native-valued-generic-point}
\lean{EnvelopeProperty.GaussCommonVerticalPoints.verticalMorphism_closedPoint}
\leanok
Put $T=S_{(\pi)}$ and let $q:\operatorname{Spec}T\to Y$ be the actual
coefficient-localization map followed by $Y.\operatorname{fromSpecStalk}(p)$.
Then $q(\mathfrak m_T)=\eta$. The entire generic-field map in the existing
original-stalk native-point construction is the actual generic map of this
SAME $q$, with every original coefficient retained.
\end{lemma}
\begin{proof}
The preceding original component-stalk theorem internally constructs
$E_x:T\simeq\mathcal O_{Y,\eta}$ and proves its equality on every
original coefficient. The structure-sheaf specialization triangle then
identifies the actual whole $q$ with
\[
 \operatorname{Spec}T\xrightarrow{\operatorname{Spec}(E_x^{-1})}
 \operatorname{Spec}\mathcal O_{Y,\eta}\longrightarrow Y.
\]
A genuine local-ring isomorphism sends the closed point to the closed
point, and the original stalk morphism sends that closed point to $\eta$.

For the generic map, compose $q$ with the actual fraction inclusion into
$\operatorname{Frac}S$. The canonical fraction-ring isomorphism between
the SAME $S$ and the SAME $Y.\mathrm{functionField}$ commutes with every
original coefficient. The actual specialization triangle from the generic
stalk to $S$ therefore identifies this composite with the existing
specified-field generic map. Equality follows on the actual coefficient
homomorphisms, rather than from a supplied field or norm comparison.
\end{proof}

\begin{theorem}[Every existing normalized original-stalk vertical point has that true center]
\label{thm:gauss-common-existing-native-vertical-center}
\uses{lem:gauss-common-original-vertical-dvr-morphism,lem:gauss-native-original-stalk-point-sequence,thm:actual-dvr-fraction-point,scheme:native-coefficient-restriction,berkovich:actual-dvr-native-model-point,berkovich:actual-dvr-native-model-center,alggeom:original-stalk-coefficient-center,lem:gauss-surface-parameter-positive-vertical-order,alggeom:original-stalk-coefficient-map,alggeom:original-stalk-coefficient-triangle,scheme:valued-point-pullback,scheme:affine-valued-point-homeomorphism,spectrum:affine-seminorm}
\lean{EnvelopeProperty.GaussCommonVerticalPoints.center_nativeVerticalPoint}
\leanok
For the SAME existing point $q_{\mathrm{vert},x}$ over the prescribed
positive Gauss radius $r<1$, its actual proper-model center is
$c(q_{\mathrm{vert},x})=\eta$. Its underlying original scheme point is
still the SAME generic point of $Y$.
\end{theorem}
\begin{proof}
The existing construction uses exactly the original DVR $T$, its actual
fraction embedding into $\operatorname{Frac}S$, and the positive scale
$-\log r/\operatorname{ord}_T(t_p)$. Positivity and actual polynomial
Gauss restriction were already derived from the original parameter.
The actual DVR fraction norm has on every original $a\in T$ exactly
the original DVR norm $\exp(-c\operatorname{ord}_T(a))$.
The existing whole native base-point restriction theorem therefore maps
its original fraction-field point to the literal original DVR base point.
Compose with the SAME $q$ from the preceding lemma. Its actual generic
map identity identifies the resulting whole native point with the
already existing $q_{\mathrm{vert},x}$.

The existing native DVR model-center theorem now gives its genuine
center as $q(\mathfrak m_T)=\eta$. The restriction to the fixed Gauss
base is an actual equality of whole coefficient valuations; changing
the proof of that restriction does not change the whole point or its
center. Generic support is the already proved underlying-point law of
the same original-stalk construction.
\end{proof}

\begin{lemma}[An actual positive Gauss point with vertical DVR center has generic support]
\label{lem:gauss-common-vertical-dvr-center-generic-support}
\uses{scheme:proper-native-center,scheme:relative-valued-points,spectrum:gauss-valuation}
\lean{EnvelopeProperty.GaussCommonVerticalPoints.genericSupport_of_verticalDVRCenter}
\leanok
Suppose $\eta$ is vertical and its actual original stalk is a DVR.
Any actual native point over the same positive Gauss radius whose true
center is $\eta$ has underlying original scheme point $\eta_Y$.
\end{lemma}
\begin{proof}
The true native support generizes its actual center. The actual range
of $Y.\operatorname{fromSpecStalk}(\eta)$ consists precisely of these
generizations. A DVR has only its zero prime and its maximal prime:
every nonzero prime is maximal by dimension one, and localness makes
it the genuine maximal ideal. The latter maps to $\eta$.

The actual global generic point also belongs to this range. It is not
$\eta$, since the generic stalk of an integral scheme is a field whereas
the actual stalk at $\eta$ is a DVR and is not a field. Consequently
the zero prime maps to the global generic point. A native support over
the prescribed positive Gauss valuation cannot be $\eta$: the original
parameter vanishes at that vertical point, but has Gauss value $r>0$.
Its preimage in the actual DVR spectrum is therefore the zero prime.
Its underlying original point is exactly the global generic point.
\end{proof}

\begin{lemma}[The same true vertical DVR center determines the whole prescribed point]
\label{lem:gauss-common-vertical-dvr-gauss-uniqueness}
\uses{lem:gauss-common-vertical-dvr-center-generic-support,lem:actual-generic-center-native-field-scale,lem:actual-generic-native-field-valuations,lem:actual-generic-center-stalk-coefficients,scheme:relative-section-coefficients,alggeom:original-stalk-coefficient-center,spectrum:gauss-valuation,alggeom:original-stalk-coefficient-map}
\lean{EnvelopeProperty.GaussCommonVerticalPoints.eq_of_same_verticalDVRCenter}
\leanok
Two actual native points over the SAME positive Gauss radius with the
SAME true vertical DVR center are the SAME whole original native point.
Their valuations agree on EVERY element of $Y.\mathrm{functionField}$.
\end{lemma}
\begin{proof}
The preceding lemma derives generic support for both actual points.
The existing genuine center classifier expresses both whole original
function-field valuations as positive real scales of the native valuation
of the SAME actual center DVR. The actual center coefficient theorem
identifies their value at the original parameter with the SAME $r$.
The parameter is nonzero and lies in the maximal ideal: its value is
strictly between zero and one, using the actual center-stalk value law.
It thus has a true positive finite native integer order $e$ in that DVR.

For scales $s_1,s_2$, the existing integer-exponential scaling formula
gives $\log r=-s_1e=-s_2e$. Since $e>0$, this forces $s_1=s_2$.
The two whole function-field valuations are therefore equal on EVERY
original fraction, including zero. The existing generic native-point
uniqueness theorem identifies the whole original native points.
No equality of norms, scales, or points is assumed.
\end{proof}

\begin{theorem}[One old original global vertex equals every existing local vertical point]
\label{thm:gauss-common-one-original-vertical-point}
\uses{berkovich:original-vertical-point-prescribed-gauss-center,thm:gauss-common-original-component-stalk,thm:gauss-common-existing-native-vertical-center,lem:gauss-common-vertical-dvr-gauss-uniqueness,alggeom:original-stalk-coefficient-injective,alggeom:original-stalk-coefficient-center,lem:gauss-common-original-component-prime,lem:actual-generic-native-field-valuations,alggeom:original-stalk-coefficient-map,lem:gauss-common-vertical-dvr-center-generic-support}
\lean{EnvelopeProperty.GaussCommonVerticalPoints.exists_commonVerticalPoint}
\leanok
Suppose the actual $f$ is dominant and $\eta$ lies in its original
zero fibre. Suppose at least one original point of $C$ has the genuine
local Noetherian prime-component data above. At the prescribed paper
radius $r=e^{-1}$, there is ONE actual original native point $q_C$ with
true center $\eta$ such that, for EVERY original $x\in C$ with actual
Noetherian stalk, prime component equation $\pi$ and
$\ker(i_x^\#)=(\pi)$, the parameter conditions $t_p\ne0$ and
$\pi\mid t_p$ are internally derived and the SAME existing normalized point satisfies
\[
 q_{\mathrm{vert},x}=q_C.
\]
Their whole original function-field valuations agree on every original
rational function. No common point, center, norm, field identification,
or local/global DVR comparison is supplied.
\end{theorem}
\begin{proof}
Use the already existing original proper dominant vertical-center
existence theorem ONCE at $\eta$ to obtain $q_C$. This is the same old
global point producer, not a new point-existence or resolution argument.
The one actual original local component equation and Noetherian stalk
internally derive that the SAME global $\eta$ stalk is a DVR by the
preceding component-stalk theorem.
For any such original $x$, the old original coefficient injectivity
theorem applied to the actual proper dominant $f$ makes $t_p\ne0$.
The actual specialization map sends $t_p$ to the original $t_\eta$,
by the actual germ-specialization square. Original verticality and the
true coefficient-center contraction put $t_\eta$ in the actual generic
stalk maximal ideal. The already proved original component-prime identity
therefore gives $t_p\in(\pi)$, hence $\pi\mid t_p$.
These original parameter facts are produced, not assumed in this final
endpoint. The actual center theorem above now derives
$c(q_{\mathrm{vert},x})=\eta$ internally.
Apply the preceding same-center uniqueness theorem to $q_C$ and this
unchanged original local point. It gives whole point equality and the
whole original function-field value equality. Thus a source convexity
hypothesis may be applied once at this ONE vertex before using the
existing arbitrary-given-unit-family profiles at different original
component points.
\end{proof}

No global SNC surface or component is constructed here. Residual leading
orders on the SAME curve, section regauging, the two principal-degree
balances and weighted degree nonnegativity remain separate producers.
